\section{Problem and decision}
\label{sec:problem}

\subsection{Supply-chain situation}
DataCo Global sells clothing, sporting goods and electronics online to 164 countries. At
checkout the customer chooses a shipping mode, and each mode carries a fixed promise: Same Day
(0 days), First Class (1), Second Class (2) or Standard Class (4). In the data, 57\% of the
orders that were not cancelled arrived late. Each late order costs goodwill credits, repeat
contacts and, eventually, customers.

\subsection{Decision, decision-maker and moment}
\begin{itemize}
  \item \textbf{Decision.} Should this order be treated as high-risk before it leaves the
        warehouse? If so, it gets priority handling or a mode upgrade, or customer service
        contacts the buyer in advance with a realistic date.
  \item \textbf{Decision-maker.} The fulfilment planner, with customer service acting on the
        flagged list.
  \item \textbf{Moment.} Order release: the order has been placed and paid, but not yet
        dispatched. At that moment the planner knows the order contents, prices and discounts,
        the customer and destination, the chosen mode and its promise, and the time of day. The
        planner does \emph{not} know the actual transit time or the final order status.
\end{itemize}

\subsection{Target and what is known at the moment of decision}
The target is \tcode{Late\_delivery\_risk} (1 = late). We verified that for every order that was
not cancelled it equals \tcode{Days for shipping (real)} $>$ \tcode{Days for shipment
(scheduled)}. Actual transit takes 0--6 days, and we assume the outcome is recorded one day after delivery,
so the label becomes known 1--7 days after the order. \Cref{sec:prep} lists the
columns we excluded because they are only known after the outcome.

\subsection{Baselines the models must beat}
The dataset does not describe DataCo's current practice, so we built baselines of increasing
strength:
\begin{enumerate}
  \item \textbf{Majority class}: treat every order as late (57\% of orders are late).
  \item \textbf{Rule A}, our proxy for current practice: flag every First and Second Class
        order, because these modes carry the tightest promises.
  \item \textbf{Rule B}: Rule A, plus Same Day orders placed at or after a cut-off hour learned
        from the training data.
  \item \textbf{Lookup table}: the historical late share for each mode, with Same Day split at
        the cut-off hour. This uses the same information as Rule B but outputs a probability.
\end{enumerate}
A model only adds value if it beats the lookup table.

\subsection{Cost of each kind of error}
\label{subsec:costs}
We assume the following costs (\Cref{tab:costs}). The dataset contains none of them, so
\Cref{sec:recommendation} tests 36 alternative combinations.

\begin{table}[!htbp]
  \centering
  \caption{Cost assumptions per order (USD). $v$ is the order value; the median is about \usd{560} in the training period.}
  \label{tab:costs}
  \small
  \begin{tabularx}{\linewidth}{l c X}
    \toprule
    \textbf{Outcome} & \textbf{Cost} & \textbf{Reasoning} \\
    \midrule
    False alarm (flagged, on time)  & \usd{10}                    & priority handling, an upgrade or a call that was not needed \\
    Missed late order               & $0.10\,v$ ($\approx\usd{56}$) & goodwill credit, repeat contacts, churn risk \\
    Caught late order               & $\usd{10} + 0.05\,v$         & action cost; half of the late-delivery cost is avoided \\
    Correctly not flagged           & 0                           & -- \\
    \bottomrule
  \end{tabularx}
\end{table}

At the median order value, a missed late order costs about five times as much as a false
alarm. For cheap orders the ratio falls below one, which drives the policy in
\Cref{sec:recommendation}.
